\usepackage[scr=rsfs,cal=euler]{mathalfa}
\multlinegap0pt

\graphicspath{{lelievre-et-al_figures}}
\hyphenation{addi-tion approx-i-mat-ed Assu-me because behav-ior decreas-ing devel-oped dimen-sion dimen-sion-al dimen-sions dyna-mics extend implies indeed intro-duced matrix meta-sta-ble oper-a-tor oper-a-tors recall result results using}
\newenvironment{enumeratei}
{\bgroup\def\theenumi{\roman{enumi}}\def\theenumii{\arabic{enumii}}\begin{enumerate}}
{\end{enumerate}\egroup}

\newcommand{\Psfrac}[2]{(\sfrac{#1}{#2})}
\newcommand{\psfrac}[2]{\sfrac{(#1)}{#2}}

\let\ts\textstyle
\let\epsilon\varepsilon

\newcommand\from{\mathchoice{\longleftarrow}{\leftarrow}{\leftarrow}{\leftarrow}}
\newcommand\mto{\mathchoice{\longmapsto}{\mapsto}{\mapsto}{\mapsto}}
\def\hto{\mathrel{\lhook\joinrel\to}}
\newcommand\To[1]{\mathchoice{\xrightarrow{\textstyle\kern4pt#1\kern3pt}}{\stackrel{#1}{\longrightarrow}}{}{}}
\newcommand\From[1]{\mathchoice{\xleftarrow{\textstyle~#1~}}{\stackrel{#1}{\longleftarrow}}{}{}}

\newcommand{\quand}{\quad\text{and}\quad}

\DeclareMathOperator{\capp}{cap}
\DeclareMathOperator{\Card}{Card}
\DeclareMathOperator{\dist}{dist}
\DeclareMathOperator{\inte}{int}
\DeclareMathOperator{\Jac}{Jac}
\DeclareMathOperator{\Span}{Span}

\usepackage{tikz}
\usetikzlibrary{arrows}
\usetikzlibrary{decorations.pathreplacing,decorations.markings}
\tikzset{arrow data/.style 2 args={%
decoration={%
markings,
mark=at position #1 with \arrow{#2}},
postaction=decorate}
}%

\newcommand{\ft}[1]{\mathsf{#1}}
\newcommand{\mc}[1]{\mathcal{#1}}
\newcommand{\mbf}[1]{\mathbf{#1}}

\let\lp\langle
\let\rp\rangle
\let\pa\partial
\let\ve\varepsilon

\DeclareMathOperator*{\argmin}{arg\,min}
\DeclareMathOperator{\diver}{div}\let\div\diver
\DeclareMathOperator{\hess}{Hess}\let\Hess\hess
\DeclareMathOperator\im{Im}\let\Im\im
\DeclareMathOperator\Ker{Ker}
\DeclareMathOperator{\range}{Ran}\let\Ran\range
\DeclareMathOperator\Sp{Sp}
\DeclareMathOperator\supp{supp}

\theoremstyle{plain}
\newtheorem{theorem}{Theorem}
\newtheorem*{theorem*}{Theorem}
\newtheorem{proposition}{Proposition}
\newtheorem{lemma}[proposition]{Lemma}
\newtheorem{corollary}[proposition]{Corollary}

\theoremstyle{definition}
\newtheorem{definition}[proposition]{Definition}
\newtheorem*{definition*}{Definition}
\newtheorem{remark}[proposition]{Remark}
\newtheorem{example}[proposition]{Example}
\newtheorem*{assumption*}{Assumption}
\newtheorem{manualassumptioninner}{Assumption}
\newenvironment{manualassumption}[1]{%
\renewcommand\themanualassumptioninner{#1}%
\manualassumptioninner}{\endmanualassumptioninner}
